% SGRMC: Sparsified Graph-Regularized Matrix Completion for Drug Repositioning
% Bioinformatics (OUP) submission
%
% Template: oup-authoring-template v1.3 (2026/05/28)
%   - modern:  Bioinformatics journal house style (required)
%   - webpdf:  hyperlink-aware PDF output
%   - large:   standard page dimensions

% ===== flushend prevention (MUST come before \documentclass) =====
\makeatletter
\expandafter\def\csname ver@flushend.sty\endcsname{9999/99/99}
\makeatother

\documentclass[webpdf,modern,large]{oup-authoring-template}

% --- OUP template metadata ---------------------------------------------------
\journaltitle{[Journal Name]}
\pubyear{[Year]}
\copyrightyear{[Year]}
\vol{[Vol]}
\issue{[Issue]}
\DOI{[DOI to be added]}

\usepackage{graphicx}
\usepackage{booktabs}
\usepackage{amsmath}
\usepackage{amssymb}
\newcounter{algline}
\usepackage{multirow}
\usepackage{xcolor}
\usepackage{hyperref}

% --- Template bug workarounds (cls v1.3) ---
\makeatletter
\providecommand{\authormark}[1]{}
\long\def\abstract#1{\let\subheading\abstsubheading\global\def\@abstract{#1}}
\let\flushend\relax
\def\societylogo{}
\makeatother


\newcommand{\R}{\mathbb{R}}
\newcommand{\E}{\mathbb{E}}
\newcommand{\tr}{\mathrm{tr}}

\begin{document}

\title{SGRMC: Sparsified Graph-Regularized Matrix Completion\\for Drug Repositioning}

\author{[Author Name(s)]}
\address{[Affiliation, Department, University, City, Country]}
\corresp{[Corresponding author: email@example.com]}

\date{}

\abstract{%
\textbf{Motivation:} Drug repositioning identifies new therapeutic uses for approved drugs at reduced cost and risk. Matrix completion methods exploit structural similarities among drugs and diseases as side information, but the similarity graphs---typically dense (all pairwise similarities are non-zero)---contain numerous weak, noisy edges that degrade regularization. We show that KNN-sparsified similarity graphs, combined with graph Laplacian regularization in a single pass, produce substantially more accurate association predictions than conventional approaches.\par
\textbf{Results:} We propose SGRMC (Sparsified Graph-Regularized Matrix Completion), a single-pass framework with three components: (1)~density-adaptive KNN graph sparsification that retains only the top-$k$ neighbors per node, removing weak similarity edges that propagate noise through the Laplacian; (2)~graph Laplacian regularization embedded directly in nuclear norm minimization; and (3)~automatic solver selection between a Neumann-approximated joint ADMM path (Plan~A, $\le 1000$ rows/cols) and a Cholesky-based two-step path (Plan~B, for larger matrices). Under 10-fold disease-centric cross-validation (CVc), SGRMC achieves AUPR of 0.1686 (Fdataset), 0.1829 (Cdataset), and 0.3724 (DNdataset), representing improvements of $+20.6$\%, $+37.9$\%, and $+26.6$\% over BNNR. Ablation experiments reveal a density-dependent division of labor: KNN sparsification drives the majority of gains on moderate-density datasets, while graph Laplacian filtering dominates on ultra-sparse data.\par
\textbf{Availability:} Source code and datasets are available at [GitHub URL: to be added upon publication].\par
\textbf{Contact:} [Corresponding author email: to be added]%
}
\maketitle

\section{Introduction}

Drug discovery costs exceed \$2.6 billion per approved compound, with development timelines spanning 10--15 years \cite{dimasi2003,paul2010}. Drug repositioning---identifying new indications for already-approved drugs---bypasses early-stage safety trials and substantially reduces both cost and risk \cite{chong2007}. Computational approaches to this problem have proliferated across multiple paradigms: network-based methods \cite{martinez2015,luo2016}, matrix factorization \cite{dai2015,ezzat2017,zheng2013}, graph neural networks \cite{fu2022,gao2023}, and matrix completion \cite{yang2019,luo2018,mongia2022}.

Matrix completion is particularly well-suited to drug repositioning because the drug--disease association matrix is partially observed and approximately low-rank---drugs with similar chemical structures tend to interact with similar disease targets, and vice versa. The nuclear norm (sum of singular values) provides the tightest convex relaxation of the rank function \cite{candes2009}, enabling tractable convex optimization via the Alternating Direction Method of Multipliers (ADMM) \cite{boyd2011}. Side information---drug--drug chemical similarity (e.g., Tanimoto coefficients over CDK descriptors \cite{steinbeck2003}) and disease--disease phenotypic similarity (e.g., MimMiner \cite{vandriel2006})---can be incorporated by augmenting the association matrix with similarity blocks, providing auxiliary structure that regularizes the completion.

However, a critical issue has received little attention: the similarity graphs used as side information are typically \emph{dense}---every drug pair has a non-zero Tanimoto similarity, every disease pair has a non-zero MimMiner score. Many of these edges are weak (e.g., Tanimoto similarity $\approx 0.1$ between structurally unrelated drugs) and contribute noise rather than signal. When such graphs are used to construct the graph Laplacian for regularization, these weak edges introduce high-frequency noise into the Laplacian eigenbasis, distorting the very manifold structure the regularizer is meant to enforce.

We propose SGRMC (Sparsified Graph-Regularized Matrix Completion), a single-pass framework built on three principles:

\begin{enumerate}
    \item \textbf{Density-adaptive KNN graph sparsification.} Before constructing the graph Laplacian, each similarity graph is sparsified to its $k$-nearest neighbors per node. The sparsification level $k$ is set adaptively based on the association matrix density: aggressive sparsification ($k=5$) for moderate-to-dense matrices where weak edges are predominantly noise; conservative sparsification ($k=50$) for ultra-sparse matrices where similarity edges are scarce information. This single preprocessing step yields the largest performance improvement observed in our experiments---$+20.6\%$ and $+37.9\%$ AUPR gains on Fdataset and Cdataset, respectively.

    \item \textbf{Graph Laplacian regularization embedded in ADMM.} We incorporate graph smoothness---the requirement that similar drugs should have similar association profiles---directly into the ADMM objective via a first-order Neumann approximation of the bilateral graph filter (Plan~A). For large matrices where per-iteration Laplacian multiplications become prohibitive, we provide a two-step path: standard completion followed by a post-hoc exact Cholesky bilateral filter (Plan~B). The solver path is selected automatically based on matrix dimensions.

    \item \textbf{Single-pass design.} The graph Laplacian is computed once from KNN-sparsified similarities and remains fixed throughout. There is no alternating refinement and no iterative graph recomputation. The completion is graph-regularized in a single pass.
\end{enumerate}

We evaluate SGRMC on three benchmark datasets spanning a 70$\times$ range of association densities (0.015\%--1.04\%) under disease-centric cross-validation (CVc), where all unlabeled entries in held-out disease columns serve as negatives. Ablation experiments on KNN sparsification level, graph filter strength, and the interaction between the two characterize the contribution of each component.

\section{Materials and Methods}

\subsection{Datasets}

We evaluate SGRMC on three benchmark datasets widely used in the drug repositioning literature \cite{yang2019,luo2016}. Each dataset consists of a drug--disease association matrix $\mathbf{A} \in \{0,1\}^{n_d \times n_s}$ (where $n_d$ is the number of drugs and $n_s$ is the number of diseases), along with drug--drug and disease--disease similarity matrices derived from chemical structures and phenotype ontologies respectively.

\begin{table*}[tb]
\centering
\caption{Dataset statistics.}
\label{tab:datasets}
\begin{tabular}{lcccc}
\toprule
Dataset & Drugs ($n_d$) & Diseases ($n_s$) & Known associations & Density \\
\midrule
Fdataset & 593 & 313 & 1,933 & $1.04 \times 10^{-2}$ \\
Cdataset & 663 & 409 & 2,532 & $9.33 \times 10^{-3}$ \\
DNdataset & 1,490 & 4,516 & 1,008 & $1.50 \times 10^{-4}$ \\
\bottomrule
\end{tabular}
\end{table*}

Drug similarities are computed as Tanimoto coefficients over CDK chemical descriptors \cite{steinbeck2003}, yielding $\mathbf{W}_{rr} \in \R^{n_d \times n_d}$. Disease similarities are computed using MimMiner \cite{vandriel2006}, which measures textual similarity between MeSH phenotype descriptions, yielding $\mathbf{W}_{dd} \in \R^{n_s \times n_s}$. Both similarity matrices take values in $[0, 1]$ with 1 on the diagonal. The three datasets span a 70$\times$ range of association densities, from moderate (Fdataset at 1.04\%) to ultra-sparse (DNdataset at 0.015\%), providing a rigorous test of robustness across density regimes.

\subsection{SGRMC: Method Description}

SGRMC integrates three components into a single-pass framework: KNN graph sparsification, graph Laplacian regularization, and nuclear norm matrix completion with ADMM. We describe each in turn before presenting the full algorithm.

\subsubsection{Preliminaries: Nuclear Norm Matrix Completion with ADMM}

Let $\mathbf{A} \in \{0,1\}^{n_d \times n_s}$ be a partially observed drug--disease association matrix, and let $\Omega = \{(i,j) : A_{ij} \text{ is observed}\}$ denote the set of known entries. The goal of matrix completion is to recover a complete matrix $\mathbf{M}$ that is low-rank and consistent with observations. The nuclear norm (sum of singular values) is the tightest convex relaxation of the rank function \cite{candes2009}, leading to:

\begin{equation}
\min_{\mathbf{X}} \|\mathbf{X}\|_* \quad \text{s.t.} \quad \mathcal{P}_\Omega(\mathbf{X}) = \mathcal{P}_\Omega(\mathbf{T}), \quad a \leq X_{ij} \leq b
\label{eq:mc}
\end{equation}

where $\mathcal{P}_\Omega$ projects onto the observed entries and $[a,b]=[0,1]$ enforces valid probability bounds. The augmented matrix $\mathbf{T}$ stacks drug--drug similarity, disease--disease similarity, and drug--disease association into a single block structure \cite{yang2019}:

\begin{equation}
\mathbf{T} = \begin{bmatrix} \mathbf{S}_{\text{drug}} & \mathbf{A}^\top \\ \mathbf{A} & \mathbf{S}_{\text{dis}} \end{bmatrix}
\label{eq:augmented}
\end{equation}

This construction encodes the intuition that similarity and association information are complementary---the side matrices $\mathbf{S}_{\text{dis}}, \mathbf{S}_{\text{drug}}$ provide auxiliary structure that regularizes the completion of $\mathbf{A}$. The optimization is solved via ADMM \cite{boyd2011}, whose W-update admits a closed-form solution and whose X-update reduces to singular value thresholding (SVT) \cite{cai2010}. Upon convergence, the completed drug--disease association matrix is extracted from the off-diagonal block of $\mathbf{X}$.

\subsubsection{Graph Laplacian Regularization}

Given a similarity graph $\mathbf{S}$ with degree matrix $\mathbf{D} = \operatorname{diag}(\mathbf{S}\mathbf{1})$, the normalized Laplacian is $\mathbf{L} = \mathbf{I} - \mathbf{D}^{-1/2} \mathbf{S} \mathbf{D}^{-1/2}$ \cite{chung1997}. The graph smoothness of a matrix $\mathbf{M}$ with respect to the drug and disease similarity graphs is measured by the Laplacian quadratic form:

\begin{equation}
\tr(\mathbf{M}^\top \mathbf{L}_{\text{dis}} \mathbf{M}) + \tr(\mathbf{M} \mathbf{L}_{\text{drug}} \mathbf{M}^\top)
\label{eq:smoothness}
\end{equation}

Minimizing this quantity encourages $\mathbf{M}$ to vary smoothly along the similarity manifold: similar drugs should have similar association profiles across diseases, and similar diseases should be associated with similar drugs. The strength of this regularization is controlled by a parameter $\alpha$, with the bilateral graph filter taking the form:

\begin{equation}
\tilde{\mathbf{M}} = (\mathbf{I} + \alpha \mathbf{L}_{\text{drug}})^{-1}
\, \mathbf{M} \,
(\mathbf{I} + \alpha \mathbf{L}_{\text{dis}})^{-1}
\label{eq:gf}
\end{equation}

Equation~\eqref{eq:gf} describes a bilateral graph low-pass filter: each row $i$ of $\mathbf{M}$ is smoothed by taking a weighted average of rows $j$ where $\mathbf{S}_{\text{drug}}(i,j)$ is large, and analogously for columns. When embedded in ADMM iterations, this jointly optimizes low-rankness and graph smoothness. When applied as a post-hoc operation, it performs spectral denoising on an already-completed matrix.

\subsubsection{KNN Graph Sparsification}

Raw similarity graphs (Tanimoto, MimMiner) are dense: every pair of drugs has a non-zero chemical similarity, and every pair of diseases has a non-zero phenotypic similarity. However, weak similarity edges between unrelated entities introduce noise into the Laplacian: the bilateral filter averages predictions across these weakly-similar neighbors, propagating errors. We therefore sparsify each similarity graph to its $k$-nearest neighbors:

\begin{equation}
\mathbf{S}_{ij}^{\text{(knn)}} =
\begin{cases}
\mathbf{S}_{ij} & \text{if } j \in \mathcal{N}_k(i) \text{ or } i \in \mathcal{N}_k(j) \\
0 & \text{otherwise}
\end{cases}
\label{eq:knn}
\end{equation}

where $\mathcal{N}_k(i)$ denotes the $k$ most similar nodes to $i$ (excluding self). The sparsified graph is symmetrized via $\mathbf{S} \gets \max(\mathbf{S}, \mathbf{S}^\top)$, and the diagonal is restored to 1. The normalized Laplacian is then computed from the sparsified graph.

The sparsification level $k$ is set adaptively based on the association matrix density $\rho$:

\begin{equation}
k = \begin{cases}
5  & \text{if } \rho > 10^{-3} \quad \text{(aggressive: weak edges are noise)} \\
50 & \text{if } \rho \le 10^{-3} \quad \text{(mild: edges are scarce information)}
\end{cases}
\label{eq:k_adaptive}
\end{equation}

On moderate-density datasets (Fdataset, Cdataset; $\rho \approx 1\%$), the dense similarity graph contains many weak edges between unrelated entities; aggressive sparsification ($k=5$) removes this noise while preserving the strongest structural relationships. On ultra-sparse matrices (DNdataset; $\rho = 0.015\%$), similarity edges are among the few sources of side information; only mild sparsification ($k=50$) is applied to avoid discarding useful signal.

\subsubsection{SGRMC: Full Formulation}

SGRMC minimizes the following objective:

\begin{multline}
\min_{\mathbf{M}}
\underbrace{\|\mathbf{M}\|_*}_{\text{low-rank}}
+ \underbrace{\frac{\alpha_{\text{mc}}}{2}\|\mathcal{P}_\Omega(\mathbf{M} - \mathbf{A})\|_F^2}_{\text{data fidelity}} \\
+ \underbrace{\alpha \cdot \big[\tr(\mathbf{M}^\top \mathbf{L}_{\text{dis}} \mathbf{M})
+ \tr(\mathbf{M} \mathbf{L}_{\text{drug}} \mathbf{M}^\top)\big]}_{\text{graph smoothness}}
\label{eq:sgrmc_obj}
\end{multline}

where $\mathbf{L}_{\text{dis}}$ and $\mathbf{L}_{\text{drug}}$ are computed once from KNN-sparsified similarity matrices (Eq.~\ref{eq:knn}) and remain fixed. The first two terms compose the standard nuclear norm completion objective; the third term, controlled by the filter strength $\alpha$, is the graph smoothness regularizer.

SGRMC provides two solver paths, selected automatically based on matrix dimensions:

\noindent\textbf{Plan A: Embedded Neumann approximation} ($\max(n_d, n_s) \le 1000$).\enspace
For small to moderate matrices, we embed the graph Laplacian regularization directly in each ADMM W-update via a first-order Neumann approximation of the bilateral graph filter:

\begin{equation}
(\mathbf{I} + \tfrac{\alpha}{10\beta}\mathbf{L})^{-1} \approx \mathbf{I} - \tfrac{\alpha}{10\beta}\mathbf{L}
\label{eq:neumann}
\end{equation}

With $\alpha = 0.7$ and $\beta = 10$, the effective per-iteration regularization strength is $\alpha/(10\beta) = 0.007$, satisfying the small-parameter requirement for first-order accuracy ($\|\tfrac{\alpha}{10\beta}\mathbf{L}\| < 1$ for any normalized Laplacian). The W-update becomes:

\begin{equation}
\mathbf{W} \gets (\mathbf{I} - \tfrac{\alpha}{10\beta}\mathbf{L}_{\text{dis}}) \cdot \mathbf{W}_{\text{off}} \cdot (\mathbf{I} - \tfrac{\alpha}{10\beta}\mathbf{L}_{\text{drug}})
\label{eq:w_update}
\end{equation}

where $\mathbf{W}_{\text{off}}$ is the off-diagonal block of the standard ADMM W-update. This joint optimization produces a solution that is simultaneously low-rank, data-consistent, and graph-smooth.

\noindent\textbf{Plan B: ADMM + exact Cholesky filter} ($\max(n_d, n_s) > 1000$).\enspace
For large matrices (e.g., DNdataset with $1490 \times 4516$), the per-iteration Laplacian multiplications in Plan~A become prohibitively expensive. We instead run standard nuclear norm completion followed by a post-hoc bilateral exact Cholesky graph filter:

\begin{equation}
\tilde{\mathbf{M}} = (\mathbf{I} + \alpha \mathbf{L}_{\text{drug}})^{-1} \cdot \mathbf{M}_{\text{completed}} \cdot (\mathbf{I} + \alpha \mathbf{L}_{\text{dis}})^{-1}
\label{eq:cholesky}
\end{equation}

The Cholesky decomposition of $(\mathbf{I} + \alpha\mathbf{L})$ is computed once per similarity matrix ($O(n^3)$ for an $n \times n$ matrix), after which each application costs $O(n^2)$ per column. This is substantially cheaper than per-iteration Laplacian multiplication in ADMM for large $n$, while producing the same graph-smoothing effect.

\subsubsection{Algorithm}

Algorithm~\ref{alg:sgrmc} summarizes the SGRMC procedure.

\begin{algorithm}[htbp]
\caption{SGRMC: Sparsified Graph-Regularized Matrix Completion}
\label{alg:sgrmc}
\small
\setlength{\tabcolsep}{0pt}
\begin{tabular}{@{\stepcounter{algline}\thealgline.\quad}p{0.82\textwidth}}
\textbf{Require:} $\mathbf{W}_{rr}, \mathbf{W}_{dd}, \mathbf{A}$ (raw similarity matrices and association matrix)\\
\textbf{Require:} $\alpha_{\text{mc}}, \beta, \alpha$ (regularization parameters)\\
\textbf{Ensure:} $\mathbf{M}$ (completed association matrix)\\
$\rho \gets |\Omega| / (n_d \cdot n_s)$ \quad // association density\\
$k \gets 5$ if $\rho > 10^{-3}$ else $50$ \quad // density-adaptive KNN sparsification\\
$\mathbf{S}_{\text{drug}} \gets \text{KNN\_sparsify}(\mathbf{W}_{rr}, k)$\\
$\mathbf{S}_{\text{dis}} \gets \text{KNN\_sparsify}(\mathbf{W}_{dd}, k)$\\
Compute $\mathbf{L}_{\text{drug}} \gets \mathcal{L}(\mathbf{S}_{\text{drug}})$, $\mathbf{L}_{\text{dis}} \gets \mathcal{L}(\mathbf{S}_{\text{dis}})$\\
Construct augmented matrix $\mathbf{T}$ via Eq.~\eqref{eq:augmented} (using $\mathbf{S}_{\text{drug}}, \mathbf{S}_{\text{dis}}$)\\
\textbf{if} $\max(n_d, n_s) \le 1000$ \textbf{then}\\
\quad $\mathbf{X} \gets \text{ADMM\_graph\_aware}(\mathbf{T}, \mathbf{L}_{\text{dis}}, \mathbf{L}_{\text{drug}}, \alpha)$ \quad // Plan A: embedded Neumann\\
\textbf{else}\\
\quad $\mathbf{X} \gets \text{ADMM\_standard}(\mathbf{T})$ \quad // standard ADMM+SVT\\
\quad $\mathbf{M} \gets \mathbf{X}[n_d:, :n_d]$ \quad // extract off-diagonal block\\
\quad $\mathbf{M} \gets (\mathbf{I} + \alpha\mathbf{L}_{\text{drug}})^{-1} \mathbf{M} (\mathbf{I} + \alpha\mathbf{L}_{\text{dis}})^{-1}$ \quad // Plan B: Cholesky filter\\
\textbf{end if}\\
$\mathbf{M}[\Omega] \gets \mathbf{A}[\Omega]$ \quad // preserve known entries\\
\textbf{return} $\mathbf{M}$
\end{tabular}
\end{algorithm}

\subsection{Evaluation Protocol}

We evaluate using 10-fold disease-centric cross-validation (CVc). For each fold, all known associations for 10\% of diseases (columns) are held out for testing, and \textbf{all} unlabeled entries in those disease columns serve as negatives. This protocol is more stringent than pair-centric cross-validation (CVa): CVc tests generalization to partially-observed diseases rather than to individual held-out associations, and evaluates against all candidate drugs rather than a subsampled negative set---yielding more conservative and clinically realistic AUPR estimates \cite{saito2015}. We report the area under the ROC curve (AUROC), area under the precision-recall curve (AUPR), and top-$K$ precision at $K=10, 20$. Given extreme class imbalance (positive rate $< 1.1\%$ in all datasets), AUPR is the primary evaluation metric.

Hyperparameters are set as follows: $\alpha_{\text{mc}} = 1$, $\beta = 10$, following standard practice \cite{yang2019}. The graph filter strength defaults to $\alpha = 0.7$ (see filter strength ablation, Section~\ref{sec:filter_ablation}). KNN sparsification uses the adaptive scheme in Eq.~\eqref{eq:k_adaptive}: $k=5$ for $\rho > 10^{-3}$, $k=50$ otherwise. ADMM convergence tolerances are $\text{tol}_1 = 2 \times 10^{-3}$, $\text{tol}_2 = 10^{-5}$ with a maximum of 300 iterations. All experiments use a fixed random seed (12345) for reproducibility.

\section{Results}
\label{sec:results}

\subsection{Main Results}

Table~\ref{tab:main_results} compares SGRMC (adaptive KNN sparsification, $\alpha = 0.7$) against BNNR, a representative nuclear norm matrix completion method \cite{yang2019}.

\begin{table*}[tb]
\centering
\caption{Main results (mean $\pm$ std, 10-fold CVc). SGRMC uses density-adaptive KNN sparsification and $\alpha = 0.7$. Best result in \textbf{bold}.}
\label{tab:main_results}
\small
\begin{tabular}{llcccc}
\toprule
Dataset & Method & AUROC & AUPR & P@10 & P@20 \\
\midrule
\multirow{2}{*}{Fdataset}
& BNNR & $0.7820 \pm 0.0369$ & $0.1398 \pm 0.0442$ & 0.43 & 0.43 \\
& \textbf{SGRMC} & $0.7867 \pm 0.0340$ & $\mathbf{0.1686 \pm 0.0538}$ & 0.51 & 0.54 \\
\midrule
\multirow{2}{*}{Cdataset}
& BNNR & $0.7943 \pm 0.0415$ & $0.1326 \pm 0.0477$ & 0.53 & 0.52 \\
& \textbf{SGRMC} & $0.8014 \pm 0.0461$ & $\mathbf{0.1829 \pm 0.0604}$ & 0.67 & 0.64 \\
\midrule
\multirow{2}{*}{DNdataset}
& BNNR & $0.8391 \pm 0.0399$ & $0.2942 \pm 0.1162$ & 0.72 & 0.59 \\
& \textbf{SGRMC} & $0.9493 \pm 0.0180$ & $\mathbf{0.3724 \pm 0.0369}$ & 0.79 & 0.70 \\
\bottomrule
\end{tabular}
\end{table*}

SGRMC substantially outperforms BNNR on all three datasets. The AUPR improvements are $+20.6\%$ (Fdataset), $+37.9\%$ (Cdataset), and $+26.6\%$ (DNdataset). AUROC improvements are most pronounced on the ultra-sparse DNdataset ($0.9493$ vs.\ $0.8391$, $+13.1\%$ relative gain), reflecting the graph Laplacian filter's ability to suppress false positives across the full ROC curve.

The improvement is not uniform across datasets. On Fdataset and Cdataset ($\rho \approx 1\%$), the gains are driven primarily by KNN sparsification (see ablation, Section~\ref{sec:knn_ablation}): sparsifying the similarity graph to $k=5$ neighbors removes noisy edges that distort the Laplacian eigenbasis. On DNdataset ($\rho = 0.015\%$), the graph Laplacian filter is the dominant contributor: the Cholesky bilateral filter denoises the completed matrix, reducing both mean prediction error and per-fold variance (AUPR std: $\pm 0.0369$ vs.\ $\pm 0.1162$ for BNNR). This variance reduction is practically significant---the bilateral filter pulls outlier predictions toward the neighborhood consensus, providing a floor on per-fold performance.

Top-$K$ precision shows a consistent pattern: SGRMC matches or exceeds BNNR on P@10 and P@20 across all datasets. The largest relative gains are on Cdataset P@10 ($0.67$ vs.\ $0.53$, $+26.4\%$), Fdataset P@20 ($0.54$ vs.\ $0.43$, $+25.6\%$), and DNdataset P@20 ($0.70$ vs.\ $0.59$, $+18.6\%$), demonstrating that the ranking quality of top candidates---those most likely to be selected for experimental validation---improves substantially.

\subsection{KNN Graph Sparsification}
\label{sec:knn_ablation}

Table~\ref{tab:knn} reports prediction performance as a function of the KNN sparsification level $k$. For each dataset, we sweep $k \in \{5, 10, 20, 50, 100\}$ with the graph filter strength fixed at $\alpha = 0.7$, and include a no-sparsification baseline ($k=0$, full dense graph) and a BNNR reference (no sparsification, no graph filter).

\begin{table*}[tb]
\centering
\caption{KNN graph sparsification ablation (mean $\pm$ std, 10-fold CVc, $\alpha=0.7$ unless noted). Best AUPR per dataset in \textbf{bold}.}
\label{tab:knn}
\small
\begin{tabular}{llcccc}
\toprule
Dataset & Configuration & AUROC & AUPR & P@10 & P@20 \\
\midrule
\multirow{7}{*}{Fdataset}
& BNNR ($k{=}0,\alpha{=}0$) & $0.7820 \pm 0.0369$ & $0.1398 \pm 0.0442$ & 0.43 & 0.43 \\
& $k=0$, $\alpha=0.7$ (dense) & $0.7767 \pm 0.0381$ & $0.1419 \pm 0.0466$ & 0.48 & 0.45 \\
& $k=5$, $\alpha=0$ (sparse only) & $0.7762 \pm 0.0361$ & $0.1515 \pm 0.0505$ & 0.49 & 0.54 \\
& \textbf{$k=5$} & $0.7867 \pm 0.0340$ & $\mathbf{0.1686 \pm 0.0538}$ & 0.51 & 0.54 \\
& $k=10$ & $0.7879 \pm 0.0302$ & $0.1665 \pm 0.0511$ & 0.57 & 0.56 \\
& $k=20$ & $0.7881 \pm 0.0353$ & $0.1521 \pm 0.0442$ & 0.51 & 0.52 \\
& $k=50$ & $0.7863 \pm 0.0368$ & $0.1416 \pm 0.0435$ & 0.48 & 0.40 \\
& $k=100$ & $0.7872 \pm 0.0381$ & $0.1439 \pm 0.0426$ & 0.47 & 0.43 \\
\midrule
\multirow{8}{*}{Cdataset}
& BNNR ($k{=}0,\alpha{=}0$) & $0.7943 \pm 0.0415$ & $0.1326 \pm 0.0477$ & 0.53 & 0.52 \\
& $k=0$, $\alpha=0.7$ (dense) & $0.7927 \pm 0.0397$ & $0.1414 \pm 0.0450$ & 0.55 & 0.56 \\
& $k=5$, $\alpha=0$ (sparse only) & $0.7888 \pm 0.0449$ & $0.1480 \pm 0.0508$ & 0.63 & 0.56 \\
& \textbf{$k=5$} & $0.8014 \pm 0.0461$ & $\mathbf{0.1829 \pm 0.0604}$ & 0.67 & 0.64 \\
& $k=10$ & $0.8036 \pm 0.0483$ & $0.1818 \pm 0.0677$ & 0.58 & 0.60 \\
& $k=20$ & $0.8047 \pm 0.0546$ & $0.1761 \pm 0.0602$ & 0.56 & 0.56 \\
& $k=50$ & $0.8031 \pm 0.0489$ & $0.1569 \pm 0.0564$ & 0.50 & 0.52 \\
& $k=100$ & $0.8008 \pm 0.0432$ & $0.1503 \pm 0.0540$ & 0.49 & 0.51 \\
\midrule
\multirow{7}{*}{DNdataset}
& BNNR ($k{=}0,\alpha{=}0$) & $0.8391 \pm 0.0399$ & $0.2942 \pm 0.1162$ & 0.72 & 0.59 \\
& $k=0$, $\alpha=0.7$ (dense) & $0.9269 \pm 0.0214$ & $0.3661 \pm 0.0320$ & 0.75 & 0.68 \\
& $k=5$, $\alpha=0$ (sparse only) & $0.8540 \pm 0.0253$ & $0.2731 \pm 0.1061$ & 0.67 & 0.59 \\
& $k=5$ & $0.8958 \pm 0.0177$ & $0.2667 \pm 0.1030$ & 0.62 & 0.57 \\
& $k=20$ & $0.9341 \pm 0.0224$ & $0.3704 \pm 0.0359$ & 0.74 & 0.64 \\
& \textbf{$k=50$} & $0.9493 \pm 0.0180$ & $\mathbf{0.3724 \pm 0.0369}$ & 0.79 & 0.70 \\
& $k=100$ & $0.9455 \pm 0.0222$ & $0.3663 \pm 0.0363$ & 0.82 & 0.74 \\
\bottomrule
\end{tabular}
\end{table*}

Three patterns emerge. First, KNN sparsification is the largest single contributor to SGRMC's performance on the moderate-density datasets. With the graph filter disabled ($\alpha=0$), sparsification alone ($k=5$) improves AUPR from 0.1398 to 0.1515 on Fdataset and from 0.1326 to 0.1480 on Cdataset---demonstrating that the Laplacian computed from a sparsified graph is a fundamentally better regularizer than one computed from the dense similarity matrix, even without explicit filtering. Adding the graph filter ($\alpha=0.7$) on top of sparsification yields further gains (0.1686 and 0.1829 respectively), showing that the two components are complementary: sparsification cleans the Laplacian eigenbasis, and the filter then leverages this cleaner basis for denoising.

Second, performance on Fdataset and Cdataset declines monotonically as $k$ increases beyond 5. At $k=100$, performance approaches the full-graph baseline, consistent with the sparsified graph converging to the dense original. This monotonic decline confirms that weak similarity edges are predominantly noise, not signal, at these densities.

Third, the pattern reverses on DNdataset: aggressive sparsification ($k=5$) degrades performance (AUPR 0.2667 vs.\ 0.3661 with the full dense graph), because similarity edges are scarce side information in this ultra-sparse regime ($\rho=0.015\%$). The optimal $k=50$ provides a modest improvement ($+1.7\%$ over the full-graph baseline). This density-dependent reversal is the motivation for our adaptive scheme (Eq.~\ref{eq:k_adaptive}).

\subsection{Graph Filter Strength}
\label{sec:filter_ablation}

Table~\ref{tab:alpha} reports the effect of varying the graph filter strength $\alpha \in \{0, 0.1, 0.3, 0.5, 0.7\}$, with KNN sparsification disabled ($k=0$, full dense graph) to isolate the effect of graph Laplacian regularization alone.

\begin{table*}[tb]
\centering
\caption{Graph filter strength ablation (mean $\pm$ std, 10-fold CVc, $k=0$). $\alpha = 0$ corresponds to BNNR without graph regularization. Best AUPR per dataset in \textbf{bold}.}
\label{tab:alpha}
\small
\begin{tabular}{llcccc}
\toprule
Dataset & $\alpha$ & AUROC & AUPR & P@10 & P@20 \\
\midrule
\multirow{5}{*}{Fdataset}
& 0.0 (BNNR) & $0.7820 \pm 0.0369$ & $0.1398 \pm 0.0442$ & 0.43 & 0.43 \\
& 0.1 & $0.7813 \pm 0.0373$ & $\mathbf{0.1422 \pm 0.0448}$ & 0.44 & 0.44 \\
& 0.3 & $0.7796 \pm 0.0378$ & $0.1407 \pm 0.0457$ & 0.44 & 0.44 \\
& 0.5 & $0.7781 \pm 0.0380$ & $0.1413 \pm 0.0466$ & 0.46 & 0.44 \\
& 0.7 & $0.7767 \pm 0.0381$ & $0.1419 \pm 0.0466$ & 0.48 & 0.45 \\
\midrule
\multirow{5}{*}{Cdataset}
& 0.0 (BNNR) & $0.7943 \pm 0.0415$ & $0.1326 \pm 0.0477$ & 0.53 & 0.52 \\
& 0.1 & $\mathbf{0.7951 \pm 0.0413}$ & $0.1340 \pm 0.0473$ & 0.52 & 0.52 \\
& 0.3 & $0.7947 \pm 0.0406$ & $0.1362 \pm 0.0474$ & 0.53 & 0.54 \\
& 0.5 & $0.7940 \pm 0.0401$ & $0.1403 \pm 0.0452$ & 0.53 & 0.54 \\
& 0.7 & $0.7927 \pm 0.0397$ & $\mathbf{0.1414 \pm 0.0450}$ & 0.55 & 0.56 \\
\midrule
\multirow{5}{*}{DNdataset}
& 0.0 (BNNR) & $0.8391 \pm 0.0399$ & $0.2942 \pm 0.1162$ & 0.72 & 0.59 \\
& 0.1 & $0.8704 \pm 0.0326$ & $0.3271 \pm 0.0965$ & 0.72 & 0.61 \\
& 0.3 & $0.9026 \pm 0.0261$ & $0.3296 \pm 0.0966$ & 0.73 & 0.64 \\
& 0.5 & $0.9181 \pm 0.0232$ & $0.3328 \pm 0.0970$ & 0.73 & 0.66 \\
& 0.7 & $\mathbf{0.9269 \pm 0.0214}$ & $\mathbf{0.3661 \pm 0.0320}$ & 0.75 & 0.68 \\
\bottomrule
\end{tabular}
\end{table*}

Graph Laplacian regularization ($\alpha > 0$) consistently improves AUPR over the no-filter baseline ($\alpha = 0$) on all three datasets, confirming that graph smoothness is a beneficial inductive bias. The benefit of stronger regularization increases with data sparsity: on Fdataset and Cdataset, the AUPR gains from $\alpha > 0$ are modest ($+1.7\%$ and $+6.6\%$ respectively at optimal $\alpha$), and the AUROC decreases slightly as $\alpha$ increases---reflecting the first-order Neumann approximation in Plan~A trading off some rank structure for graph smoothness. On DNdataset (Plan~B, exact Cholesky), both AUROC and AUPR improve monotonically with $\alpha$, with no sign of saturation at $\alpha = 0.7$. This suggests that on extremely sparse matrices, stronger filtering yields monotonically better denoising.

The DNdataset AUPR standard deviation notably decreases from $\pm 0.1162$ ($\alpha=0$) to $\pm 0.0320$ ($\alpha=0.7$): the bilateral filter reduces not only mean prediction error but also per-fold variance, pulling outlier predictions toward the neighborhood consensus.

We fix $\alpha = 0.7$ as the default for SGRMC. While Fdataset achieves a marginally higher AUPR at $\alpha = 0.1$ (0.1422 vs.\ 0.1419, a difference of 0.0003), this gap is far smaller than the fold-to-fold variance and is eliminated once KNN sparsification is added (Table~\ref{tab:knn}, where $\alpha=0.7$ with $k=5$ achieves the best result). A single fixed $\alpha$ simplifies the method and avoids dataset-specific tuning.

\subsection{Computational Cost}

Table~\ref{tab:cost} summarizes per-fold runtime.

\begin{table}[htbp]
\centering
\caption{Per-fold runtime (minutes, 10-fold CVc).}
\label{tab:cost}
\begin{tabular}{lccc}
\toprule
Method & Fdataset & Cdataset & DNdataset \\
\midrule
BNNR & ${<}1$ & ${<}1$ & ${\sim}2$ \\
SGRMC & ${<}1$ & ${<}1$ & ${\sim}4$ \\
\bottomrule
\end{tabular}
\end{table}

On Fdataset and Cdataset (matrices $\le 1000$, Plan~A), SGRMC adds negligible overhead: the per-iteration Neumann approximation (two matrix multiplications per W-update) is cheap for matrices of this size. On DNdataset ($1490 \times 4516$, Plan~B), SGRMC adds approximately 2 minutes per fold, attributable to the Cholesky decomposition and bilateral filter application. This overhead is modest relative to the $+26.6\%$ AUPR gain. KNN sparsification itself adds negligible cost ($O(n^2 \log k)$ per similarity matrix for a naive implementation, or $O(n^2)$ with partitioning). The automatic solver selection (Plan~A vs.\ Plan~B) ensures that SGRMC remains computationally tractable across dataset sizes.

\section{Discussion}

\subsection{Why Graph Sparsification Helps}

The effectiveness of KNN sparsification can be understood through the spectral properties of the graph Laplacian. A dense similarity graph contains many edges with small positive weights (e.g., Tanimoto similarities of 0.1--0.3 between structurally dissimilar drugs). In the normalized Laplacian $\mathbf{L} = \mathbf{I} - \mathbf{D}^{-1/2}\mathbf{S}\mathbf{D}^{-1/2}$, these weak edges contribute small off-diagonal entries that, when accumulated across all $\sim n$ neighbors, introduce high-frequency noise into the graph Fourier basis. The bilateral filter (Eq.~\ref{eq:gf}) projects predictions onto the leading eigenvectors of $(\mathbf{I}+\alpha\mathbf{L})^{-1}$, which approximate the low-frequency (smooth) components of the graph. When the Laplacian is contaminated by weak-edge noise, these leading eigenvectors are distorted away from the true community structure.

Sparsifying to $k$-nearest neighbors removes this noise: edges below the $k$-th largest similarity per node are zeroed, eliminating their contribution to the Laplacian. The resulting leading eigenvectors better reflect the true block structure of the similarity graph (clusters of chemically related drugs, biologically related diseases). When the bilateral filter projects the completed matrix onto these cleaner eigenvectors, it preserves genuine community signal while more effectively attenuating prediction noise.

This mechanism explains the complementary relationship between sparsification and filtering observed in Tables~\ref{tab:knn} and~\ref{tab:alpha}. On Fdataset and Cdataset, sparsification alone ($k=5, \alpha=0$) already improves AUPR by $+8.4\%$ and $+11.6\%$ over BNNR, because a cleaner Laplacian improves the augmented matrix conditioning even without explicit filtering. Adding the graph filter on top of a sparsified graph yields further gains ($+20.6\%$ and $+37.9\%$ total), because the filter now operates on a clean eigenbasis where genuine community structure dominates.

\subsection{Density-Dependent Behavior}

The optimal sparsification level $k$ depends critically on association matrix density, motivating our adaptive scheme (Eq.~\ref{eq:k_adaptive}). This dependence reflects a fundamental trade-off between noise removal and information preservation.

On moderate-density datasets ($\rho \approx 1\%$), the similarity graphs contain $\sim n^2/2$ edges each---far more than the number of known associations (${\sim}2000$). In this regime, weak similarity edges are predominantly noise: they connect entities that are not meaningfully related, and their aggregate effect on the Laplacian eigenbasis is detrimental. Aggressive sparsification ($k=5$) removes these edges and retains only the strongest structural relationships, which tend to correspond to genuine chemical or phenotypic similarity.

On the ultra-sparse DNdataset ($\rho = 0.015\%$, 1,008 known associations among 6.7 million possible pairs), the situation is reversed. The association signal is so scarce that any structural information---even weak similarity edges---may carry useful signal. Aggressive sparsification ($k=5$) discards too much of this information, degrading performance below the no-sparsification baseline. Mild sparsification ($k=50$) provides a small but consistent improvement, suggesting that even in this regime, the very weakest edges (those below the 50th percentile per node) are still noise.

This density-dependent reversal has practical implications: sparsification strategies should be calibrated to the data regime. Our adaptive threshold ($\rho = 10^{-3}$, Eq.~\ref{eq:k_adaptive}) provides a simple heuristic; a continuous function of density (e.g., $k \propto 1/\sqrt{\rho}$) would be a natural extension for datasets with intermediate density.

\subsection{Design Choices and Rationale}

SGRMC embodies three deliberate design choices:

\noindent\textbf{KNN-sparsified graphs over dense similarity graphs.}
Dense similarity graphs contain weak edges that introduce noise into the Laplacian eigenbasis. KNN sparsification removes these edges before regularization, yielding a cleaner graph Fourier basis. This preprocessing step accounts for the majority of SGRMC's performance gain on moderate-density datasets and has negligible computational cost.

\noindent\textbf{Single-pass over alternating refinement.}
When the similarity graph is fixed and not iteratively recomputed, there is no graph to refine---the completion is graph-regularized in a single pass. SGRMC's single-pass design eliminates the computational overhead of alternating schemes while preserving the core mechanism of graph-regularized completion.

\noindent\textbf{Fixed filter strength with adaptive sparsification.}
The graph filter strength $\alpha = 0.7$ is fixed across all datasets, while only the sparsification level $k$ adapts to association density. The filter strength ablation (Section~\ref{sec:filter_ablation}) shows that $\alpha = 0.7$ is near-optimal across all three datasets, and the KNN ablation (Section~\ref{sec:knn_ablation}) shows that adaptive $k$ covers the density-dependent sparsification needs. This design separates two concerns that are often conflated: how strongly to enforce smoothness (fixed) and which edges to consider when measuring smoothness (adaptive).

\subsection{Comparison with Prior Work}

The HN-DRES benchmark \cite{li2024hn_dres} provides a comprehensive recent evaluation of 28 network-based drug repositioning methods under pair-centric cross-validation (CVa). Direct per-dataset comparison between SGRMC and these methods is not currently feasible: the benchmark reports results under CVa (with 1:1 negative subsampling), while our evaluation uses CVc (all-unlabeled-as-negative). Under CVa, AUPR values are systematically inflated (typically 0.90--0.98) because only a small, balanced subset of negatives is evaluated; under CVc, AUPR values drop to 0.13--0.37, reflecting the realistic challenge of distinguishing true associations among a vast candidate pool.

The two protocols answer different questions. CVa asks: ``can the method separate known associations from a random sample of unlabeled pairs?'' CVc asks: ``can the method rank true associations above all possible candidates for a disease with only partial observed associations?'' The latter is more clinically relevant for drug repositioning, where the practical task is to prioritize a small number of candidates from the entire drug space for experimental validation. We encourage future benchmarking efforts to adopt CVc as a standardized protocol alongside CVa.

\subsection{Limitations and Future Work}

Several limitations merit discussion. First, SGRMC currently uses two similarity modalities (chemical and phenotypic). Multi-view extensions---incorporating gene expression, protein--protein interactions, or drug side-effect profiles---are compatible with the framework: the graph Laplacian regularization naturally extends to multiple graphs per entity type via multi-kernel Laplacian combinations.

Second, the Plan~A Neumann approximation (Eq.~\ref{eq:neumann}) is accurate only when the effective filter strength $\alpha/(10\beta)$ is small relative to $1/\|\mathbf{L}\|$. While this holds for our parameter settings ($\alpha/(10\beta) = 0.007$), substantially higher filter strengths would require second-order approximations or exact Cholesky solves even for small matrices. The Plan~B Cholesky path does not have this limitation.

Third, the adaptive KNN threshold ($\rho = 10^{-3}$) is set between Cdataset ($\rho = 9.33 \times 10^{-3}$) and DNdataset ($\rho = 1.50 \times 10^{-4}$), and validated on these two regimes. Datasets with intermediate density ($10^{-3}$--$10^{-2}$) may benefit from a continuous sparsification function rather than a hard threshold, which we leave to future work.

Fourth, while SGRMC demonstrates strong performance under CVc, prospective experimental validation of top-ranked predictions would strengthen translational relevance.

\section{Conclusion}

We introduced SGRMC, a single-pass framework for drug repositioning that combines density-adaptive KNN graph sparsification with graph Laplacian regularization in nuclear norm matrix completion. Three principles underpin the method: (1)~sparsify similarity graphs to their $k$-nearest neighbors before constructing the Laplacian, with $k$ set adaptively based on association density; (2)~enforce graph smoothness via Laplacian regularization embedded in ADMM (Plan~A, for small matrices) or applied as a post-hoc Cholesky bilateral filter (Plan~B, for large matrices); and (3)~complete the matrix in a single pass without alternating refinement.

Under disease-centric cross-validation (CVc), SGRMC achieves AUPR of 0.1686 (Fdataset), 0.1829 (Cdataset), and 0.3724 (DNdataset)---representing relative improvements of $+20.6\%$, $+37.9\%$, and $+26.6\%$ over BNNR. Ablation experiments reveal a density-dependent division of labor: KNN sparsification is the dominant performance driver on moderate-density datasets, where it removes noisy edges that distort the Laplacian eigenbasis; graph Laplacian filtering dominates on the ultra-sparse dataset, where it provides critical spectral denoising.

Our findings challenge a widely held assumption in graph-regularized matrix completion---that denser similarity graphs provide better regularization. Under realistic, all-unlabeled-as-negative evaluation, the opposite is true: sparsified graphs, by excluding weak and noisy similarity edges, produce a cleaner manifold structure and substantially more accurate predictions. This principle---sparsify the structural prior before regularizing the completion---may generalize beyond drug repositioning to other matrix completion tasks where side information graphs are available but noisy.

\section*{Data Availability}

The datasets used in this study are publicly available from prior publications \cite{yang2019,luo2016}. Source code for SGRMC and all experiments is available at [GitHub URL: to be added upon publication].

\section*{Funding}

This work was supported by [Funding: grant number(s) to be added].

\section*{Conflict of Interest}

None declared.

\section*{Author Contributions}

[Author initials] conceived the method, implemented the algorithms, conducted experiments, analyzed results, and wrote the manuscript.

\section*{Acknowledgments}

We thank the authors of prior matrix completion benchmarks \cite{yang2019} for making their code and datasets publicly available.

\bibliographystyle{oup-abbrvnat}
\bibliography{references}

\end{document}
